% PAPER A (submission track), forked 2026-08-17 per the external referee's structure
% recommendation and the principal's order. The crisp k=2 paper: model, balance lemma,
% two-defect theorem, rigidity classification, closing three-defect teaser. The full
% three-defect/uniform-program material lives in ../main.tex (the living ladder draft).
% Source of truth for every theorem: the campaign wiki entries named in comments.
% Gates before submission: principal line-read + external cold-read referee.
% TODO(acquire): Aigner 1985 deposited copy for the width-two equality citation.
% TODO(principal): authorship.
\documentclass[11pt]{article}
\usepackage{amsmath,amssymb,amsthm}
\usepackage{booktabs}
\usepackage[margin=1.2in]{geometry}

\newtheorem{theorem}{Theorem}[section]
\newtheorem{lemma}[theorem]{Lemma}
\newtheorem{proposition}[theorem]{Proposition}
\newtheorem{corollary}[theorem]{Corollary}
\theoremstyle{definition}
\newtheorem{definition}[theorem]{Definition}
\theoremstyle{remark}
\newtheorem{remark}[theorem]{Remark}

\newcommand{\pr}{\operatorname{Pr}}
\DeclareMathOperator{\bal}{bal}

\title{The $1/3$--$2/3$ Conjecture for Posets\\ at Chain-Deletion Distance Two}
\author{}% TODO(principal): authorship
\date{Draft of August 16, 2026}

\begin{document}
\maketitle

\begin{abstract}
A finite poset is a \emph{$k$-defect chain} if it becomes a chain after the deletion of
$k$ elements. We prove the $1/3$--$2/3$ conjecture for all two-defect chains: every such
poset that is not itself a chain has an incomparable pair $(u,v)$ with
$1/3\le\pr(u\prec v)\le 2/3$. The class contains posets of width three and does not bound the number of
incomparabilities per element, and we exhibit members escaping each previously proved
class known to us. We prove equality rigidity: every ordinal-sum-indecomposable poset of
defect exactly two has a strictly balanced eligible count, and among
chain-plus-two-incomparable-defect presentations the unique non-strict member is the
three-element poset $Q$ (a two-element chain plus an isolated point, itself of defect
one); this extends Aigner's width-two equality classification to an infinite family of
width-three posets, whose non-attainment of $1/3$ the order-$14$ census could previously
only observe empirically. A closing section takes the first step onto the three-defect ladder: an
infinite three-defect family attaining the boundary exactly, and the precise place where
the two-defect arithmetic stops.
\end{abstract}

\section{Introduction}
% Positioning per the banked novelty sweep (wiki: operator-novelty-sweep-two-defect-
% theorem-new-rigidity-width). All class attributions verified against produced sources.
The $1/3$--$2/3$ conjecture --- posed independently by Kislitsyn \cite{kislitsyn},
Fredman \cite{fredman}, and Linial \cite{linial} in connection with sorting under partial
information --- asserts that every finite poset $P$ that is not a chain contains an
incomparable pair $(u,v)$ whose orientation probability under the uniform distribution on
linear extensions satisfies $1/3\le\pr(u\prec v)\le 2/3$. Kahn and Saks \cite{kahn-saks}
proved the first uniform positive bound; the strongest general bound is due to Brightwell,
Felsner, and Trotter \cite{bft}. The conjecture is known for partial orders of width two
(Linial \cite{linial}, sharpened by Sah \cite{sah}), height two \cite{tgf}, posets with at
most $14$ elements \cite{deloof,gupta}, posets in which every element is incomparable to
at most six others (Peczarski \cite{peczarski-thin}), series-parallel posets, $N$-free
posets and polytrees (Zaguia \cite{zaguia-nfree,zaguia-polytree}), semiorders (Brightwell
\cite{brightwell-semi}), several lattice and Young-diagram families (Olson and Sagan
\cite{olson-sagan}), and is preserved by lexicographic sums \cite{dolores-cuenca}. Chan
and Pak \cite{chan-pak} survey the area.

Every finite poset becomes a chain after deleting some number of elements; call the least
such number its \emph{defect}. The conjecture restricted to defect $k$ interpolates, as
$k$ grows, from the merging problems of the classical literature to the full conjecture.
Defect one is contained in width two. This paper settles defect two completely, classifies
its extremal cases, and closes with the first rung of the three-defect ladder.

Three subfamilies of the two-defect class were previously known and are not claimed here:
members whose two defect windows coincide carry a defect-exchanging automorphism, a case
settled by Peczarski's automorphism theorem \cite{peczarski-aut} (and, for cycle-free
members, already by Ganter, H\"afner, and Poguntke \cite{ghp}, who obtain a $1/2$-balanced
pair); members of width two are Linial--Sah; members whose windows have
length at most two are semiorders. The class as a whole --- with windows of unbounded length, hence unboundedly many
incomparabilities at a single element, and with genuine width-three members --- is
contained in neither the width-two class nor any bounded-incomparability class, and its
non-containment in the remaining listed classes is witnessed explicitly. With $m=2$ and
windows $I=[0,1]$, $J=[1,2]$, the four elements $(x,c_1,c_2,y)$ induce an $N$ ($x<c_2$,
$c_1<c_2$, $c_1<y$, all other pairs incomparable), both as an induced suborder and in
the Hasse diagram --- so the class is contained in neither the series-parallel nor the
$N$-free class. With $m=3$ and windows $I=[1,2]$, $J=[1,3]$, the Hasse diagram contains
the cycle $c_1\,x\,c_3\,c_2$ --- so it is not contained in the forest-cover class.
Unbounded windows exclude semiorders, and height excludes the height-two class. The incomparable-defect
members are precisely the interval orders with exactly two nontrivial intervals; general
interval orders remain open.

Section~2 constructs the insertion-window model and its state bijection; Section~3
proves the balance lemma; Section~4 assembles the Two-Defect Chain Theorem; Section~5
proves the equality rigidity classification; Section~6 poses the next problem. Throughout, $e(P)$ is the number of linear
extensions of $P$, $\pr(u\prec v)$ the fraction orienting $u$ before $v$, and
$\#L(P;\,\cdot)$ the number of extensions with the displayed property; windows are
integer intervals of insertion gaps against the fixed chain $C$.

\section{The doubled-diagonal model}
% Source: wiki doubled-diagonal-overlapping-interval-balance-decision (round 5, sol),
% independently re-derived in round 6 (supervisor-verdict-work-d4e5c56a7791).
Let $C=(c_1<\cdots<c_m)$ be a chain and let $x,y$ be two further elements with
$x\parallel y$. Consistency with the chain forces the admissible insertion gaps of $x$ to
form a window $I=[a,b]\subseteq\{0,\dots,m\}$ (gap $i$ meaning that exactly $i$ chain
elements precede $x$): the relations are $c_t<x$ for $t\le a$ and $x<c_{t}$ for $t>b$.
Similarly $y$ carries a window $J=[c,d]$.

\begin{lemma}[Overlap]\label{lem:overlap}
If $x\parallel y$ then $I\cap J\neq\emptyset$. Conversely, for any two nonempty
overlapping windows $I,J$, imposing the chain order, the threshold relations of both
windows, and no relation between $x$ and $y$ defines a valid poset in which $x\parallel y$
and no hidden comparabilities arise.
\end{lemma}

\begin{proof}
If $b<c$ then $x<c_{b+1}\le c_c<y$ by transitivity, contradicting $x\parallel y$; the case
$d<a$ is symmetric. For the converse, every pair of gaps $(i,j)\in I\times J$ is realized
by merging $x$ and $y$ into the chain at those gaps, and at a common gap both relative
orders of $x,y$ occur, so $x\parallel y$ and the prescribed windows are exactly attained.
% imported from: round-6 verdict items "Insertion-Gap Structural Model" and "Overlap
% Converse for Two Insertion Intervals"
\end{proof}

\begin{lemma}[State bijection and counts]\label{lem:states}
Let $p=|I|$, $q=|J|$, $K=I\cap J$, $h=|K|$. The linear extensions of $P$ are in bijection
with the states $(i,j)\in I\times J$ with $i\ne j$, together with two ordered states for
each common gap $i=j\in K$; hence
\[
N=e(P)=(pq-h)+2h=pq+h.
\]
The orientation counts of the incomparable pairs are
\begin{align*}
D&=\#\{(i,j)\in I\times J: i<j\}+h &&\text{for the pair }(x,y),\\
A_t&=q\,\#\{i\in I: i<t\}+\#\{i\in K: i<t\} &&\text{for }(x,c_t),\ a<t\le b,\\
B_t&=p\,\#\{j\in J: j<t\}+\#\{j\in K: j<t\} &&\text{for }(y,c_t),\ c<t\le d,
\end{align*}
and exactly the displayed thresholds correspond to incomparable pairs; outside them the
orientation is forced.
\end{lemma}

\begin{proof}
A linear extension is determined by the pair of gaps, except that when $x$ and $y$ occupy
a common gap the two relative orders give two extensions --- the doubled diagonal. For
$D$: every off-diagonal state with $i<j$ contributes, and exactly the $x$-before-$y$
diagonal state contributes, giving the displayed count. For $A_t$ with $a<t\le b$: $x<c_t$
exactly when the $x$-gap satisfies $i<t$; each such $i$ pairs with $q$ choices of $j$,
plus one extra state when $i\in K$ (the doubled diagonal), and symmetrically for $B_t$.
Incomparability of the displayed pairs is witnessed by explicit gap choices realizing both
orientations.
% imported from: r5 artifact section 1 (state model), verified against the m<=4
% topological-sort audit (259 cases, 0 failures) recorded in the same artifact
\end{proof}

\section{The balance lemma}
% Source: same r5/r6 entries; full derivation banked verbatim in the r5 artifact.
\begin{theorem}[Overlapping Two-Interval Balance Lemma]\label{thm:balance}
For every pair of nonempty overlapping windows, at least one of $D$, an eligible $A_t$,
or an eligible $B_t$ lies in $[N/3,\,2N/3]$.
\end{theorem}

The proof rests on one discrete intermediate-value fact and exact jump accounting.

\begin{lemma}[Discrete crossing lemma]\label{lem:crossing}
Let $x_1<\cdots<x_r$ be integers with $x_1\le N/3$, $x_r\ge 2N/3$, and
$x_{k+1}-x_k\le N/3$ for all $k$. Then some $x_k\in[N/3,\,2N/3]$.
\end{lemma}

\begin{proof}
Take the first $k$ with $x_k\ge N/3$. If $k=1$ then $x_1=N/3$ exactly and we are done;
otherwise $x_{k-1}<N/3$, so $x_k=x_{k-1}+(x_k-x_{k-1})<N/3+N/3=2N/3$.
\end{proof}

Exchanging the defect names replaces $D$ by $N-D$ and swaps the $A$ and $B$ families;
translation and order reversal also preserve the assertion. After these symmetries exactly
two canonical forms remain: \emph{crossing} ($a\le c\le b\le d$; then $h=p-s$ where $s=c$
after translating $a$ to $0$, and $q\ge h$) and \emph{nested} ($J\subseteq I$; then
$h=q$).

\paragraph{The crossing form.} Direct summation gives
\[
N=pq+h,\qquad D=pq-\tfrac{h(h-1)}{2},\qquad
A_t=qt+\max(0,t-s),\qquad B_{s+u}=pu+\min(u,h),
\]
with exact consecutive jumps $A_{t+1}-A_t\in\{q,q+1\}$ and
$B_{s+u+1}-B_{s+u}\in\{p,p+1\}$. If $N\ge 3(q+1)$ then $p\ge2$, the $A$-sequence runs
from $A_1\le q+1\le N/3$ to $A_{p-1}=N-(q+1)\ge 2N/3$ with jumps at most $q+1\le N/3$,
and Lemma~\ref{lem:crossing} produces a balanced eligible $A_t$. Symmetrically, if
$N\ge3(p+1)$ the $B$-sequence does the same. If neither large-case condition holds, the
inequalities $q(p-3)<3-h$ and $p(q-3)<3-h$ with $1\le h\le\min(p,q)$ force exactly fifteen
triples $(p,q,h)$, each of which carries an explicit balanced witness:

\begin{center}
\begin{tabular}{@{}lrrl@{}}\toprule
$(p,q,h)$ & $s$ & $N$ & witness \\ \midrule
$(1,1,1)$ & 0 & 2 & $D=1$\\
$(1,2,1)$ & 0 & 3 & $D=2$\\
$(1,3,1)$ & 0 & 4 & $B_{s+1}=2$\\
$(1,4,1)$ & 0 & 5 & $B_{s+1}=2$\\
$(2,1,1)$ & 1 & 3 & $D=2$\\
$(2,2,1)$ & 1 & 5 & $A_1=2$\\
$(2,2,2)$ & 0 & 6 & $D=3$\\
$(2,3,1)$ & 1 & 7 & $A_1=3$\\
$(2,3,2)$ & 0 & 8 & $D=5$\\
$(3,1,1)$ & 2 & 4 & $A_2=2$\\
$(3,2,1)$ & 2 & 7 & $A_2=4$\\
$(3,2,2)$ & 1 & 8 & $D=5$\\
$(3,3,1)$ & 2 & 10 & $A_2=6$\\
$(3,3,2)$ & 1 & 11 & $A_2=7$\\
$(4,1,1)$ & 3 & 5 & $A_2=2$\\ \bottomrule
\end{tabular}
\end{center}

\paragraph{The nested form.} Here $N=(p+1)q$. If $q\ge2$, take $u=\lceil q/3\rceil$: then
$1\le u\le q-1$ and $q\le 3u\le 2q$, so the eligible count $B_{s+u}=(p+1)u$ satisfies
$N\le 3B_{s+u}\le 2N$ \emph{exactly} --- the factorization of $N$ makes the nested case
rigid. If $q=1$ then $N=p+1$, $D=s+1$, and $A_t=t+\mathbf1_{\{s<t\}}$; for $p\ge5$ the
$A$-sequence starts at most $2$, ends at least $p-1$, and jumps by $1$ or $2$, so
Lemma~\ref{lem:crossing} applies, while $p\le4$ is a finite table of explicit witnesses.
% imported from: r5 artifact sections 4-5; the closed interval is used throughout and no
% continuity is invoked (r5 artifact section 6). Corroborating audits (r5/r6): all 5,551
% overlapping endpoint tuples at m=12; adversarial canonical scan of 4,040,100 states
% p,q<=200; zero failures. Audits are checks, not premises.

\begin{remark}[$D$ alone does not suffice]
The count $D$ alone is not always balanced: at $(p,q,h)=(2,2,1)$ one has $N=5$ and
$D=4>2N/3$, while the eligible $A_1=2$ is balanced. Any proof must use the chain-cut
counts. % source: r6 verdict, "D-Only Balance Route Refuted"
\end{remark}

\section{The two-defect chain theorem}
\begin{theorem}[Two-Defect Chain Theorem]\label{thm:twodefect}
Every finite non-chain poset $P$ containing a chain $C$ with $|P\setminus C|=2$ has an
incomparable pair $(u,v)$ with $1/3\le\pr(u\prec v)\le 2/3$.
\end{theorem}

\begin{proof}
Write $P\setminus C=\{x,y\}$. If $x$ and $y$ are comparable, then $C\cup\{x<y\}$ is a
two-chain cover, so $P$ has width at most two and the theorem of Linial \cite{linial}
(in the sharpened form of Sah \cite{sah}) applies. If $x\parallel y$, their windows
overlap by Lemma~\ref{lem:overlap}; Lemma~\ref{lem:states} identifies the orientation
counts of the genuine incomparable pairs, and Theorem~\ref{thm:balance} produces one in
the middle third.
\end{proof}

\section{The extremal classification}
% Sources: r7 (supervisor verdict "DD-EQ Rigidity Proved") + r8 independent re-derivation
% (supervisor-verdict-work-5fbb947628b0). ATTRIBUTION: the width-two part is Aigner
% \cite{aigner}: width-two posets with balance constant exactly 1/3 are the ordinal sums
% of singletons and copies of Q. The width-three content is new; the order-14 census
% \cite{gupta} observes it empirically and notes it is beyond Aigner's argument.
Write $Q$ for the three-element poset with exactly one relation. Aigner \cite{aigner}
proved that the width-two posets with balance constant exactly $1/3$ are precisely the
ordinal sums of singletons and copies of $Q$. The census of Gupta \cite{gupta}
determines the extremal balance data through order $14$: exactly $128$ classes attain
$1/3$, every one an ordinal sum of singletons and copies of $Q$ --- in particular all of
width two. The following theorem proves this non-attainment for the full (infinite,
width-three-containing) class. Note the distinction between the \emph{presentation}
(a chain plus two designated defects, the model of Section~2, which any poset of defect
at most two admits) and the \emph{exact defect} $n-h(P)$: the extremal member $Q$
arises as a one-element chain plus two defects but has exact defect one.

\begin{theorem}[Equality rigidity]\label{thm:rigidity}
\emph{(a)} Every ordinal-sum-indecomposable poset of defect exactly two has an eligible
orientation count strictly inside the open interval $(N/3,\,2N/3)$; in particular its
balance constant exceeds $1/3$.
\emph{(b)} Among ordinal-sum-indecomposable posets presented as a chain plus two
incomparable defects, the unique member all of whose eligible counts avoid the open
interval is $Q$ --- which has defect one. Every other member has an eligible count
strictly inside.
\end{theorem}

The proof classifies the \emph{simultaneous avoiders} --- parameter tuples for which every
eligible count $C$ satisfies $3C\le N$ or $3C\ge2N$ --- and intersects them with an
indecomposability criterion:

\begin{proposition}[Avoider classification]\label{prop:avoiders}
In the canonical crossing form, simultaneous avoidance holds exactly for
$(p,q,h,s)=(1,2,1,0)$ and $(2,1,1,1)$. In the nested form, exactly for
$(p,q,s)=(2,1,0)$ and $(2,1,1)$.
\end{proposition}

\begin{proposition}[Indecomposability criterion]\label{prop:indec}
The two-defect model with windows $I=[a,b]$, $J=[c,d]$ is ordinal-sum-indecomposable iff
its incomparability graph is connected, iff every $t\in\{1,\dots,m\}$ satisfies $a<t\le b$
or $c<t\le d$; under the overlap hypothesis this is equivalent to $\min(a,c)=0$ and
$\max(b,d)=m$.
\end{proposition}

\begin{proof}[Proof of Proposition~\ref{prop:avoiders}]
% transcribed from paper/derivations/r7_capture_02.md (r7, sol), independently
% re-derived in r8 (r8_capture_01/02, luna) — two derivations on record.
Write $\mathrm{AV}(C)$ for the condition $3C\le N$ or $3C\ge2N$. We first sharpen
Lemma~\ref{lem:crossing} under avoidance: if $x_1\le N/3$, $x_r\ge2N/3$, all jumps are at
most $N/3$, and \emph{every} term satisfies $\mathrm{AV}$, then the first term
$\ge N/3$ equals $N/3$ exactly, and its successor equals $2N/3$ exactly (the predecessor
of the first crossing term is strictly below $N/3$, so the jump bound puts that term
strictly below $2N/3$; avoidance then forces equality at $N/3$, and the next term can
neither be in the open middle nor overshoot $2N/3$).

\emph{Large regimes.} In the crossing form the $A$-sequence has $A_1\le q+1$,
$A_{p-1}=N-(q+1)$, and jumps in $\{q,q+1\}$. If $N\ge3(q+1)$, all-$A$ avoidance forces
$N/3\in\{q,q+1\}$ by the sharpened lemma, hence $N=3(q+1)$, i.e.\ $q(p-3)=3-h$. With
$1\le h\le\min(p,q)$ the solutions are $(p,q,h)=(3,q,3)$ for $q\ge3$, $(4,2,1)$, and
$(5,1,1)$; the first is killed at $q=3$ by $D=6$, $N=12$ (strict middle) and for $q>3$
by the symmetric $B$-argument (which would force $N=3(p+1)=12$, impossible), while the
last two carry the strict witnesses $A_2=4$, $N=9$ and $A_3=3$, $N=6$. Symmetrically,
$N\ge3(p+1)$ with all-$B$ avoidance forces $p(q-3)=3-h$, whose solutions
$(p,3,3)_{p\ge3}$, $(2,4,1)$, $(1,5,1)$ are killed the same way. Hence every survivor
satisfies $N<3(q+1)$ and $N<3(p+1)$.

\emph{Residual list.} The strict inequalities $q(p-3)<3-h$ and $p(q-3)<3-h$ force
$p,q\le4$ and exactly the fifteen triples of the table in
Section~\ref{thm:balance}. All but two carry an eligible count strictly inside the open
middle third:
\[
\begin{array}{llllll}
(1,1,1)\!: D{=}1,N{=}2 & (1,3,1)\!: B_1{=}2,N{=}4 & (1,4,1)\!: B_1{=}2,N{=}5\\
(2,2,1)\!: A_1{=}2,N{=}5 & (2,2,2)\!: D{=}3,N{=}6 & (2,3,1)\!: A_1{=}3,N{=}7\\
(2,3,2)\!: D{=}5,N{=}8 & (3,1,1)\!: A_2{=}2,N{=}4 & (3,2,1)\!: A_2{=}4,N{=}7\\
(3,2,2)\!: D{=}5,N{=}8 & (3,3,1)\!: A_2{=}6,N{=}10 & (3,3,2)\!: A_2{=}7,N{=}11\\
(4,1,1)\!: A_2{=}2,N{=}5
\end{array}
\]
(each displayed value $C$ satisfies $N<3C<2N$). The remaining two triples avoid: at
$(1,2,1)$, $N=3$ with $D=B_1=2$, both at the boundary $3C=2N$; at $(2,1,1)$, $N=3$ with
$D=2$ and $A_1=1$, at the boundaries $3C=2N$ and $3C=N$. Note $D-N/2=(pq-h^2)/2\ge0$
throughout, so $D$ never realizes the low branch in this orientation.

\emph{Nested form.} Here $N=(p+1)q$ and $B_{s+u}=(p+1)u$. If $q\ge2$ and $3\nmid q$,
then $u=\lceil q/3\rceil$ gives $q<3u<2q$, a strict-middle $B$; if $q=3k\ge6$, take
$u=k+1$. If $q=3$: $p=3$ gives $D=6$, $N=12$ strict middle, and for $p\ge4$ the
$A$-sequence starts at most $4<N/3=p+1$, ends at least $3p-1>2N/3$, with jumps at most
$4$, so Lemma~\ref{lem:crossing} gives a strict term. For $q=1$: $N=p+1$ and
$\{D\}\cup\{A_t\}=\{1,\dots,p\}$; $p=1$ fails at $D=1$, $N=2$, $p\ge3$ contains an
integer strictly inside, and $p=2$ survives for $s\in\{0,1\}$ with values $\{1,2\}$ at
the boundaries of $N=3$.
\end{proof}

\begin{proof}[Proof of Proposition~\ref{prop:indec}]
The incomparability graph has the edge $xy$, no edges among chain vertices, and $c_t$
adjacent to $x$ exactly for $a<t\le b$, to $y$ exactly for $c<t\le d$. After exchange
assume $a\le c$; overlap gives $c\le b$, so the union of the two adjacency windows is the
single index interval $(a,\max(b,d)]$, and the graph is connected exactly when that
interval is all of $\{1,\dots,m\}$ --- that is, $\min(a,c)=0$ and $\max(b,d)=m$.
Connectedness of the incomparability graph is equivalent to ordinal-sum
indecomposability, since the canonical ordinal-sum factors are precisely the
incomparability components. % banked: canonical-ordinal-sum-factorization theorem
\end{proof}

Intersecting the avoider tuples (with their translations, defect exchanges, and order
reversals) with the criterion leaves exactly the four interval presentations
$(L,F),(F,L),(F,R),(R,F)$, where $L=[0,0]$, $R=[1,1]$, $F=[0,1]$ at $m=1$: endpoint
coverage forces $m=s+q-1$, collapsing both survivor tuples to one-element chains with
the displayed windows, each isomorphic to $Q$ (extensions $abc,acb,cab$; pair counts
$(2,1)$ and $(1,2)$). Every other arithmetic survivor has a nontrivial chain prefix or
suffix and is a chain--$Q$--chain ordinal sum, hence decomposable. This proves
part~(b) of Theorem~\ref{thm:rigidity}.

Part (a) follows. Let $P$ be indecomposable of defect exactly two, and let $C$ be a
maximum chain, $P\setminus C=\{x,y\}$. If $x$ and $y$ are comparable, $C\cup\{x<y\}$ is
a two-chain cover, so $P$ has width two; Linial \cite{linial} gives $\delta(P)\ge1/3$,
and by Aigner's classification \cite{aigner} equality holds only for ordinal sums of
singletons and copies of $Q$ --- all decomposable except $Q$ itself, whose defect is
one. Hence $\delta(P)>1/3$: some incomparable pair lies strictly inside. If $x\parallel
y$, part (b) applies to the presentation, and its unique avoider $Q$ has defect one,
not two; so $P$ carries a strict eligible count.

\begin{remark}
An eligible count with $N/3<C<2N/3$ yields a strict witness; if all eligible counts avoid
the open interval, Theorem~\ref{thm:balance} forces equality at $N/3$ or $2N/3$ --- which
is the situation classified above. A corroborating exact scan through $p,q\le60$ (73{,}810
canonical cases) found no further avoiders. % r7; audits are checks, not premises
\end{remark}

\section{The next rung: three defects}
% Sources: r9/r10 banked aligned theorem (two independent derivations); the k=3 boundary
% from the r13 bridge verdict. The full three-defect program lives in the companion
% living draft (../main.tex in the campaign record).
The machinery of this paper is the $k=2$ instance of a general weighted-gap model: for
$k$ defects, linear extensions correspond to feasible gap assignments weighted by the
extension counts of the tied blocks. The first three-defect result is already in reach of
a symmetry argument, and it lands exactly on the conjecture's boundary.

\begin{theorem}[Aligned $Q$-defect chains]
Let $P\setminus C=\{x,y,z\}$ induce $x<y$ with $z$ incomparable to both, all three
insertion windows equal to a common interval $I$ with $p=|I|\ge2$. Then
$N=e(P)=p(p+1)(p+2)/2$ and
\[
\pr(x\prec z)=\pr(z\prec y)=2/3\qquad\text{exactly.}
\]
\end{theorem}

\begin{proof}[Proof sketch]
States are triples of gaps $(i,j,k)$ with $i\le j$, weighted $1$, $2$, or $3$ by how many
defects share a gap. Delete the relation $x<y$: the three defects become interchangeable
and relabeling splits the extensions into six classes of equal size $\binom{p+2}{3}$.
Restoring $x<y$ retains exactly the classes $xyz$, $xzy$, $zxy$; two of the three have
$x\prec z$ and two have $z\prec y$.
\end{proof}

Thus the aligned family generalizes the extremal poset $Q$ of
Theorem~\ref{thm:rigidity} to every window size, with its defect pairs pinned to the
boundary. Beyond the aligned chamber the two-defect arithmetic genuinely stops: with
three defects, tied fibers are no longer pair-additive (a fully tied gap carries the
extension count of the induced three-element defect order --- $6$, $3$, $2$, or $1$ ---
rather than a product of pairwise corrections), and orientation counts become asymmetric
between the defect pairs. Deciding the full three-defect analogue of
Theorem~\ref{thm:twodefect}, and the extremal classification above width three, are the
natural next problems; the interval-order connection (incomparable-defect $k$-defect
chains are interval orders with $k$ nontrivial intervals) suggests the ladder climbs
toward genuinely open territory.

\begin{thebibliography}{99}
\bibitem{kislitsyn} S.~S. Kislitsyn, \emph{Finite partially ordered sets and their
associated sets of permutations}, Matematicheskie Zametki \textbf{4} (1968), 511--518.
\bibitem{fredman} M.~Fredman, \emph{How good is the information theory bound in sorting?},
Theoretical Computer Science \textbf{1} (1976), 355--361.
\bibitem{linial} N.~Linial, \emph{The information-theoretic bound is good for merging},
SIAM J. Computing \textbf{13} (1984), 795--801.
\bibitem{kahn-saks} J.~Kahn and M.~Saks, \emph{Balancing poset extensions}, Order
\textbf{1} (1984), 113--126.
\bibitem{aigner} M.~Aigner, \emph{A note on merging}, Order \textbf{2} (1985), 257--264.
\bibitem{ghp} B.~Ganter, G.~H\"afner, and W.~Poguntke, \emph{On linear extensions of
ordered sets with a symmetry}, Discrete Mathematics \textbf{63} (1987), 153--156.
% VERIFIED 2026-08-18 (web, produced source): GHP = Discrete Math 63 (1987) 153-156,
% cycle-free + nontrivial automorphism -> 1/2 pair; Peczarski 2006 = GPC (hence 1/3-2/3)
% for any poset with a nontrivial automorphism. Both cited accordingly in the intro.
\bibitem{brightwell-semi} G.~Brightwell, \emph{Semiorders and the 1/3--2/3 conjecture},
Order \textbf{5} (1989), 369--380.
\bibitem{tgf} W.~T. Trotter, W.~G. Gehrlein, and P.~C. Fishburn, \emph{Balance theorems
for height-2 posets}, Order \textbf{9} (1992), 43--53.
\bibitem{bft} G.~Brightwell, S.~Felsner, and W.~T. Trotter, \emph{Balancing pairs and the
cross product conjecture}, Order \textbf{12} (1995), 327--349.
\bibitem{peczarski-thin} M.~Peczarski, \emph{The Gold Partition Conjecture for 6-thin
posets}, Order \textbf{25} (2008), 91--103.
\bibitem{peczarski-aut} M.~Peczarski, \emph{The gold partition conjecture}, Order
\textbf{23} (2006), 89--95. % TODO(verify): exact home of the automorphism theorem
\bibitem{deloof} K.~De Loof, B.~De Baets, and H.~De Meyer, \emph{Counting linear
extension majority cycles in partially ordered sets on up to 13 elements}, Computers \&
Mathematics with Applications \textbf{59} (2010), 1541--1547.
\bibitem{zaguia-nfree} I.~Zaguia, \emph{The 1/3--2/3 conjecture for $N$-free ordered
sets}, Electronic J. Combinatorics \textbf{19}(2) (2012), P29.
\bibitem{zaguia-polytree} I.~Zaguia, \emph{The 1/3--2/3 conjecture for ordered sets whose
cover graph is a forest}, Order \textbf{36} (2019), 335--347.
\bibitem{olson-sagan} E.~J. Olson and B.~E. Sagan, \emph{On the 1/3--2/3 conjecture},
Order \textbf{35} (2018), 581--596.
\bibitem{sah} A.~Sah, \emph{Improving the $\frac13$--$\frac23$ conjecture for width two
posets}, Combinatorica \textbf{41} (2021), 99--126.
\bibitem{dolores-cuenca} E.~Dolores-Cuenca, A.~Guzm\'an-S\'aenz, and S.~Kim, \emph{The
gold partition conjecture and the lexicographic sum of posets}, arXiv:2410.12494.
\bibitem{chan-pak} S.~H. Chan and I.~Pak, survey. % TODO(verify): exact reference
\bibitem{gupta} R.~Gupta, \emph{Balance constants, majority cycles, and the Gold
Partition Conjecture through fourteen elements}, arXiv:2607.23926 (2026).
\end{thebibliography}

\end{document}
